\newpage
\section{Dùng công cụ Excel, SPSS, Python và R tính toán}

\subsection{Phương trình hồi quy tuyến tính}

\subsubsection{Sử dụng Excel}

Dựa trên dữ liệu ở Bảng \ref{fig:table1}, ta thực hiện phân tích hồi quy tuyến tính bội
với biến phụ thuộc là $Y$ (Output of Product) và hai biến độc lập là
$X_1$ (Productivity) và $X_2$ (Produced Cost) với mức ý nghĩa $\alpha = 5\%$.

Kết quả phân tích hồi quy bằng Excel được trình bày trong các hình
\ref{fig:reg_stat}, \ref{fig:anova}, và \ref{fig:coef}.

\begin{figure}[H]
\centering
\includegraphics[width=0.3\linewidth]{LAB03/pics/regression_statistics.png}
\caption{Regression Statistics từ kết quả Excel}
\label{fig:reg_stat}
\end{figure}

\begin{figure}[H]
\centering
\includegraphics[width=0.8\linewidth]{LAB03/pics/anova_table.png}
\caption{Bảng ANOVA của mô hình hồi quy}
\label{fig:anova}
\end{figure}

\begin{figure}[H]
\centering
\includegraphics[width=01\linewidth]{LAB03/pics/regression_coefficients.png}
\caption{Hệ số hồi quy của các biến độc lập}
\label{fig:coef}
\end{figure}

Từ kết quả hồi quy, ta thu được phương trình hồi quy tuyến tính bội như sau:

\begin{equation}
Y = 2.4748 + 3.1875X_1 + 0.2848X_2
\end{equation}

Trong đó:

\begin{itemize}
\item $Y$: Output of Product
\item $X_1$: Productivity
\item $X_2$: Produced Cost
\end{itemize}

\subsubsection*{Giải thích các chỉ số hồi quy}

\textbf{Multiple R}

Multiple R = 0.8309 thể hiện hệ số tương quan giữa giá trị dự đoán và
giá trị thực tế của biến phụ thuộc. Giá trị này gần 1 cho thấy mối quan hệ
khá mạnh giữa các biến độc lập và biến phụ thuộc.

\textbf{R-square}

Giá trị $R^2 = 0.6904$ cho thấy khoảng 69.04\% sự biến động của biến
Output of Product được giải thích bởi hai biến độc lập $X_1$ và $X_2$.

\textbf{Adjusted R-square}

Adjusted $R^2 = 0.6492$ là hệ số xác định đã được điều chỉnh theo số lượng
biến độc lập trong mô hình, cho thấy mô hình vẫn giữ được mức độ giải thích
tốt đối với dữ liệu.

\textbf{Significance F}

Giá trị Significance F = 0.0001515 nhỏ hơn mức ý nghĩa $\alpha = 0.05$,
do đó mô hình hồi quy là có ý nghĩa thống kê.

\textbf{Coefficient of Independent Variables}

Hệ số của các biến độc lập cho biết mức độ thay đổi của $Y$
khi các biến độc lập thay đổi một đơn vị.

\begin{itemize}
\item Khi Productivity ($X_1$) tăng 1 đơn vị, Output of Product tăng
khoảng 3.1875 đơn vị (giữ các yếu tố khác không đổi).
\item Khi Produced Cost ($X_2$) tăng 1 đơn vị, Output of Product tăng
khoảng 0.2848 đơn vị.
\end{itemize}

\textbf{P-value}

Giá trị P-value của hai biến độc lập đều nhỏ hơn 0.05,
do đó cả hai biến Productivity và Produced Cost đều có ảnh hưởng
có ý nghĩa thống kê đến Output of Product.

\subsubsection*{Kết luận}

Kết quả phân tích hồi quy cho thấy mô hình hồi quy tuyến tính bội
có ý nghĩa thống kê. Hai biến độc lập Productivity và Produced Cost
đều có ảnh hưởng đáng kể đến Output of Product, trong đó
Productivity có tác động mạnh hơn so với Produced Cost.

\subsubsection{Phương trình hồi quy tuyến tính bằng R}

Dựa trên dữ liệu trong Bảng \ref{fig:table1}, ta xây dựng mô hình hồi quy
tuyến tính bội với biến phụ thuộc là $Y$ (Output of Product) và hai biến độc lập
là $X_1$ (Productivity) và $X_2$ (Produced Cost). Phân tích được thực hiện bằng
ngôn ngữ R.

\subsubsubsection*{Code R}

\begin{verbatim}
# Create dataset
X1 <- c(2.5,3.4,3.5,3.6,3.4,3.3,3.6,3.5,2.5,
        3.1,3.2,3.1,3.5,3.5,3.4,3.2,2.5,2.7)

X2 <- c(30,30,45,40,40,45,40,40,30,
        45,40,45,45,40,30,30,30,30)

Y <- c(19.3,22.4,24.5,26.9,25.3,22.7,25.6,27.2,21.6,
       25.6,24,24.3,27.3,27.5,21.2,16.7,18.4,19.6)

data <- data.frame(X1,X2,Y)

# Multiple linear regression
model <- lm(Y ~ X1 + X2, data=data)

summary(model)
\end{verbatim}

\subsubsubsection*{Kết quả hồi quy trong R}

Kết quả hồi quy tuyến tính từ R được trình bày trong Hình
\ref{fig:r_regression_output}.

\begin{figure}[H]
\centering
\includegraphics[width=0.85\linewidth]{LAB03/pics/R_regression_output.png}
\caption{Kết quả hồi quy tuyến tính bội trong R}
\label{fig:r_regression_output}
\end{figure}

\subsubsubsection*{Phương trình hồi quy}

Từ kết quả phân tích hồi quy, ta thu được phương trình hồi quy tuyến tính bội:

\begin{equation}
Y = 2.4748 + 3.1875X_1 + 0.2848X_2
\end{equation}

Trong đó:

\begin{itemize}
\item $Y$: Output of Product
\item $X_1$: Productivity
\item $X_2$: Produced Cost
\end{itemize}

\subsubsubsection*{Giải thích các chỉ số hồi quy}

\textbf{Multiple R}

Multiple R thể hiện hệ số tương quan giữa giá trị dự đoán và giá trị thực tế
của biến phụ thuộc. Giá trị này gần 1 cho thấy mối quan hệ mạnh giữa
các biến độc lập và biến phụ thuộc.

\textbf{R-square}

Giá trị $R^2 = 0.6904$ cho thấy khoảng 69.04\% sự biến động của
Output of Product được giải thích bởi hai biến độc lập
Productivity và Produced Cost.

\textbf{Adjusted R-square}

Adjusted $R^2 = 0.6492$ điều chỉnh hệ số xác định theo số lượng biến
độc lập trong mô hình, giúp đánh giá chính xác hơn độ phù hợp của mô hình.

\textbf{Significance F}

Giá trị Significance F nhỏ hơn mức ý nghĩa $0.05$ cho thấy
mô hình hồi quy có ý nghĩa thống kê.

\textbf{Coefficient of Independent Variables}

\begin{itemize}
\item Khi Productivity ($X_1$) tăng 1 đơn vị, Output of Product tăng
khoảng 3.1875 đơn vị (giữ các yếu tố khác không đổi).
\item Khi Produced Cost ($X_2$) tăng 1 đơn vị, Output of Product tăng
khoảng 0.2848 đơn vị.
\end{itemize}

\textbf{P-value}

P-value của hai biến độc lập đều nhỏ hơn 0.05,
cho thấy cả hai biến Productivity và Produced Cost đều có
ảnh hưởng có ý nghĩa thống kê đến Output of Product.

\subsubsubsection*{Kết luận}

Kết quả phân tích hồi quy cho thấy mô hình hồi quy tuyến tính bội
có ý nghĩa thống kê và giải thích được khoảng 69\% sự biến động
của Output of Product. Trong đó, Productivity có tác động mạnh hơn
so với Produced Cost.


\subsubsection{Sử dụng SPSS}

Để xây dựng mô hình hồi quy tuyến tính bội, dữ liệu trong Bảng \ref{fig:table1}
được nhập vào phần mềm SPSS với ba biến: Productivity ($X_1$),
Produced Cost ($X_2$) và Output of Product ($Y$).

Trong SPSS, ta thực hiện các bước sau:

\begin{itemize}
\item Chọn \textbf{Analyze}
\item Chọn \textbf{Regression}
\item Chọn \textbf{Linear}
\item Đặt \textbf{Output of Product (Y)} vào ô \textbf{Dependent}
\item Đặt \textbf{Productivity (X1)} và \textbf{Produced Cost (X2)} vào ô \textbf{Independent}
\item Nhấn \textbf{OK} để chạy mô hình hồi quy
\end{itemize}

\subsubsection*{Kết quả hồi quy}

Kết quả phân tích hồi quy trong SPSS được trình bày trong các
Hình \ref{fig:spss_model_summary}, \ref{fig:spss_anova}
và \ref{fig:spss_coefficients}.

\begin{figure}[H]
\centering
\includegraphics[width=0.7\linewidth]{LAB03/pics/SPSS_model_summary.png}
\caption{Model Summary trong SPSS}
\label{fig:spss_model_summary}
\end{figure}

\begin{figure}[H]
\centering
\includegraphics[width=0.8\linewidth]{LAB03/pics/SPSS_anova.png}
\caption{Bảng ANOVA của mô hình hồi quy trong SPSS}
\label{fig:spss_anova}
\end{figure}

\begin{figure}[H]
\centering
\includegraphics[width=0.9\linewidth]{LAB03/pics/SPSS_coefficients.png}
\caption{Bảng hệ số hồi quy trong SPSS}
\label{fig:spss_coefficients}
\end{figure}

\subsubsection*{Phương trình hồi quy}

Từ kết quả hồi quy, ta thu được phương trình hồi quy tuyến tính bội:

\begin{equation}
Y = 2.4748 + 3.1875X_1 + 0.2848X_2
\end{equation}

Trong đó:

\begin{itemize}
\item $Y$: Output of Product
\item $X_1$: Productivity
\item $X_2$: Produced Cost
\end{itemize}

\subsubsection*{Giải thích các chỉ số hồi quy}

\textbf{Multiple R}

Multiple R thể hiện hệ số tương quan giữa giá trị dự đoán và
giá trị thực tế của biến phụ thuộc.

\textbf{R-square}

Giá trị $R^2 = 0.6904$ cho thấy khoảng 69.04\% sự biến động
của Output of Product được giải thích bởi hai biến độc lập
Productivity và Produced Cost.

\textbf{Adjusted R-square}

Adjusted $R^2 = 0.6492$ điều chỉnh hệ số xác định theo số lượng
biến độc lập trong mô hình.

\textbf{Significance F}

Giá trị Significance F = 0.0001515 nhỏ hơn mức ý nghĩa
$\alpha = 0.05$, cho thấy mô hình hồi quy có ý nghĩa thống kê.

\textbf{Coefficient of Independent Variables}

\begin{itemize}
\item Khi Productivity ($X_1$) tăng 1 đơn vị, Output of Product
tăng khoảng 3.1875 đơn vị.
\item Khi Produced Cost ($X_2$) tăng 1 đơn vị, Output of Product
tăng khoảng 0.2848 đơn vị.
\end{itemize}

\textbf{P-value}

P-value của hai biến độc lập đều nhỏ hơn 0.05,
cho thấy cả hai biến đều có ảnh hưởng có ý nghĩa
thống kê đến Output of Product.

\subsubsection*{Kết luận}

Kết quả phân tích hồi quy cho thấy mô hình hồi quy tuyến tính
bội có ý nghĩa thống kê. Hai biến độc lập Productivity và
Produced Cost đều có ảnh hưởng đáng kể đến Output of Product.


\subsubsection{Sử dụng Python}

Phân tích hồi quy tuyến tính bội cũng được thực hiện bằng Python
với thư viện \texttt{statsmodels}. Biến phụ thuộc là $Y$ (Output of Product),
và hai biến độc lập là $X_1$ (Productivity) và $X_2$ (Produced Cost).

\subsubsection*{Python Code}

\begin{verbatim}
import pandas as pd
import statsmodels.api as sm

data = {
"X1":[2.5,3.4,3.5,3.6,3.4,3.3,3.6,3.5,2.5,
      3.1,3.2,3.1,3.5,3.5,3.4,3.2,2.5,2.7],
"X2":[30,30,45,40,40,45,40,40,30,
      45,40,45,45,40,30,30,30,30],
"Y":[19.3,22.4,24.5,26.9,25.3,22.7,25.6,27.2,21.6,
     25.6,24,24.3,27.3,27.5,21.2,16.7,18.4,19.6]
}

df = pd.DataFrame(data)

X = df[['X1','X2']]
X = sm.add_constant(X)
Y = df['Y']

model = sm.OLS(Y, X).fit()

print(model.summary())

b0 = model.params['const']
b1 = model.params['X1']
b2 = model.params['X2']

print("\nRegression Equation:")
print(f"Y = {b0:.4f} + {b1:.4f}X1 + {b2:.4f}X2")
\end{verbatim}

\subsubsection*{Kết quả hồi quy}

\begin{figure}[H]
\centering
\includegraphics[width=1\linewidth]{LAB03/pics/python_regression_output.png}
\caption{Kết quả hồi quy tuyến tính bội bằng Python}
\label{fig:python_regression}
\end{figure}

\subsubsection*{Phương trình hồi quy}

Từ kết quả hồi quy, ta thu được phương trình hồi quy tuyến tính bội:

\begin{equation}
Y = 2.4748 + 3.1875X_1 + 0.2848X_2
\end{equation}

\subsubsection*{Giải thích các chỉ số}

\textbf{Multiple R}

Multiple R được tính bằng căn bậc hai của $R^2$:

\[
R = \sqrt{0.6904} \approx 0.8309
\]

\textbf{R-square}

\[
R^2 = 0.6904
\]

Khoảng 69.04\% sự biến động của Output of Product được giải thích
bởi hai biến độc lập.

\textbf{Adjusted R-square}

Adjusted $R^2 = 0.6492$ cho thấy mô hình vẫn phù hợp sau khi
điều chỉnh theo số lượng biến độc lập.

\textbf{Significance F}

F-statistic = 16.73 với p-value = 0.0001515, nhỏ hơn mức ý nghĩa
0.05 nên mô hình hồi quy có ý nghĩa thống kê.

\textbf{Coefficient và P-value}

\begin{itemize}
\item Productivity ($X_1$): hệ số 3.1875, p-value = 0.049
\item Produced Cost ($X_2$): hệ số 0.2848, p-value = 0.006
\end{itemize}

Cả hai biến đều có ảnh hưởng có ý nghĩa thống kê đến Output of Product.

\subsection{Nonlinear Regression using Excel}

Mô hình hồi quy phi tuyến được xây dựng dựa trên phương trình:

\[
Y = b_0 + b_1X_1 + b_2X_2 + b_3\frac{1}{X_1} + b_4\frac{1}{X_1X_2}
\]

Hai biến mới được tạo trong Excel:

\[
Z_1 = \frac{1}{X_1}, \quad Z_2 = \frac{1}{X_1X_2}
\]

Sau đó thực hiện Regression với các biến:

\[
X_1, X_2, Z_1, Z_2
\]

\begin{figure}[H]
\centering
\includegraphics[width=0.9\linewidth]{LAB03/pics/excel_nonlinear_regression.png}
\caption{Kết quả hồi quy phi tuyến trong Excel}
\end{figure}

Từ kết quả hồi quy, ta thu được phương trình:

\[
Y = -183.8573 + 41.6464X_1 - 1.2767X_2 + 580.1149\frac{1}{X_1}
- 7048.6430\frac{1}{X_1X_2}
\]

Giá trị $R^2 = 0.8407$ cho thấy mô hình phi tuyến giải thích
84.07\% sự biến động của Output of Product.

So sánh với mô hình tuyến tính ($R^2 = 0.690$),
mô hình phi tuyến có độ phù hợp cao hơn.

\subsubsection{Nonlinear Regression using R}

Mô hình hồi quy phi tuyến được xây dựng bằng cách tạo hai biến mới:

\[
Z_1 = \frac{1}{X_1}, \qquad Z_2 = \frac{1}{X_1X_2}
\]

Sau đó thực hiện hồi quy:

\[
Y = b_0 + b_1X_1 + b_2X_2 + b_3Z_1 + b_4Z_2
\]

\begin{figure}[H]
\centering
\includegraphics[width=0.9\linewidth]{LAB03/pics/R_nonlinear_regression.png}
\caption{Kết quả hồi quy phi tuyến trong R}
\end{figure}

Từ kết quả phân tích hồi quy, ta thu được phương trình:

\[
Y = -183.8573 + 41.6464X_1 - 1.2767X_2 + 580.1149\frac{1}{X_1}
- 7048.6430\frac{1}{X_1X_2}
\]

Giá trị $R^2 = 0.8407$ cho thấy mô hình phi tuyến giải thích
84.07\% sự biến động của Output of Product.

\subsubsection{Nonlinear Regression using SPSS}

Trong SPSS, hai biến mới được tạo bằng chức năng
\textit{Transform → Compute Variable}:

\[
Z_1 = \frac{1}{X_1}, \qquad Z_2 = \frac{1}{X_1X_2}
\]

Sau đó thực hiện hồi quy tuyến tính với các biến:

\[
X_1, X_2, Z_1, Z_2
\]

\begin{figure}[H]
\centering
\includegraphics[width=0.9\linewidth]{LAB03/pics/SPSS_nonlinear_regression.png}
\caption{Kết quả hồi quy phi tuyến trong SPSS}
\end{figure}

Từ kết quả phân tích hồi quy, ta thu được phương trình:

\[
Y = -183.857 + 41.646X_1 - 1.277X_2 + 580.115\frac{1}{X_1}
- 7048.643\frac{1}{X_1X_2}
\]

Giá trị $R^2 = 0.8407$ cho thấy mô hình phi tuyến
giải thích khoảng 84\% sự biến động của Output of Product.

\subsubsection{Nonlinear Regression using Python}

Mô hình hồi quy phi tuyến được xây dựng bằng Python bằng cách
tạo thêm hai biến:

\[
Z_1 = \frac{1}{X_1}, \qquad Z_2 = \frac{1}{X_1X_2}
\]

Sau đó thực hiện hồi quy:

\[
Y = b_0 + b_1X_1 + b_2X_2 + b_3Z_1 + b_4Z_2
\]

\begin{figure}[H]
\centering
\includegraphics[width=0.9\linewidth]{LAB03/pics/python_nonlinear_regression.png}
\caption{Kết quả hồi quy phi tuyến trong Python}
\end{figure}

Từ kết quả phân tích, ta thu được phương trình:

\[
Y = -183.8573 + 41.6464X_1 - 1.2767X_2 + 580.1149\frac{1}{X_1}
- 7048.6430\frac{1}{X_1X_2}
\]

Giá trị $R^2 = 0.8407$ cho thấy mô hình phi tuyến giải thích
84\% sự biến động của Output of Product.